% ---------------------------------------------------------------------------------
% 4. CONCLUSIONS
% Guide: largest supported effects; scratch vs fine-tuning; difficulties, limitations,
% realistic next steps.
% ---------------------------------------------------------------------------------
\section{Conclusions}\label{sec:conclusions}

\begin{itemize}
  \item \textbf{Largest effects.} CUSTOM: GAP head, more width and a longer schedule (0.44 $\to$ 0.50). FT: the ImageNet pretraining itself, with a small backbone learning rate (0.546 in 30 epochs).
  \item \textbf{Scratch vs fine-tuning.} FT is clearly better on validation, but both score about the same on the public subset (0.49 vs 0.48).
  \item \textbf{Main limitation.} The intermediate classes are hard even on training images, and reusing one validation set for selection inflates its scores ($\sim$0.06). The original FireRisk paper reports only 61--65\% accuracy with $\sim$91k images \citep{shen2023firerisk}.
  \item \textbf{Next steps.} $k$-fold cross-validation for honest selection, losses that use the class order, remote-sensing pretraining and higher input resolution.
\end{itemize}
